\chapter{Physical Constants, Units, and Notation}
\label{app:notation}

This appendix collects the physical constants, unit conventions, and notational
choices used throughout the monograph. It complements the Nomenclature in the
front matter.

\section{Physical constants}
\begin{longtable}{@{}p{6.8cm}p{2.8cm}p{4.4cm}@{}}
\toprule
\textbf{Constant} & \textbf{Symbol} & \textbf{Value (SI)} \\
\midrule
\endhead
Elementary charge & $e$ & \SI{1.602176634e-19}{\coulomb} \\
Electron mass & $m_{\mathrm{e}}$ & \SI{9.1093837015e-31}{\kilogram} \\
Proton mass & $m_{\mathrm{p}}$ & \SI{1.67262192369e-27}{\kilogram} \\
Boltzmann constant & $k_{\mathrm{B}}$ & \SI{1.380649e-23}{\joule\per\kelvin} \\
Vacuum permittivity & $\varepsilon_{0}$ & \SI{8.8541878128e-12}{\farad\per\metre} \\
Speed of light & $c$ & \SI{2.99792458e8}{\metre\per\second} \\
Planck constant & $h$ & \SI{6.62607015e-34}{\joule\second} \\
Atomic mass unit & $u$ & \SI{1.66053906660e-27}{\kilogram} \\
Avogadro constant & $N_{\mathrm{A}}$ & \SI{6.02214076e23}{\per\mole} \\
\bottomrule
\end{longtable}

\section{Unit conventions}
Plasma temperatures are expressed as energies in electronvolts
(\SI{1}{\electronvolt} $\leftrightarrow$ \SI{11604}{\kelvin} to four significant
figures). Number densities are quoted in \si{\per\cubic\metre}; where the
low-temperature-plasma literature is followed, values may also appear in
\si{\per\cubic\centi\metre} (\SI{1}{\per\cubic\centi\metre} $=$
\SI{1e6}{\per\cubic\metre}). Energy fluxes are given in \si{\watt\per\square\metre}
and particle fluxes in \si{\per\square\metre\per\second}. Pressures follow the SI
pascal, with the equivalence \SI{1}{\pascal} $\approx$ \SI{7.5}{mTorr} noted where
processing pressures are discussed.

\section{Frequently used relations}
\begin{longtable}{@{}p{5.2cm}p{9.2cm}@{}}
\toprule
\textbf{Quantity} & \textbf{Expression} \\
\midrule
\endhead
Debye length & $\lambda_{\mathrm{D}} = \sqrt{\varepsilon_{0} k_{\mathrm{B}} T_{\mathrm{e}} / (n_{\mathrm{e}} e^{2})}$ \\
Electron plasma frequency & $\omega_{\mathrm{pe}} = \sqrt{n_{\mathrm{e}} e^{2} / (\varepsilon_{0} m_{\mathrm{e}})}$ \\
Bohm velocity & $u_{\mathrm{B}} = \sqrt{k_{\mathrm{B}} T_{\mathrm{e}} / m_{\mathrm{i}}}$ \\
Ion (Bohm) flux & $\Gamma_{\mathrm{i}} = n_{\mathrm{s}} u_{\mathrm{B}}$ \\
Floating-sheath potential drop & $\Delta\phi \approx \dfrac{k_{\mathrm{B}} T_{\mathrm{e}}}{2e}\,\ln\!\left(\dfrac{m_{\mathrm{i}}}{2\pi m_{\mathrm{e}}}\right)$ \\
Ion energy flux & $\Gamma_{E} = \displaystyle\int \tfrac{1}{2} m_{\mathrm{i}} v^{2}\, v\, f_{\mathrm{i}}(v)\, \mathrm{d}v$ \\
\bottomrule
\end{longtable}

\section{Notational conventions}
Vectors are set in bold ($\bm{E}$, $\bm{v}$); scalars in italic. Subscripts
$\mathrm{e}$ and $\mathrm{i}$ denote electron and ion quantities; a subscript
$\mathrm{s}$ denotes evaluation at the sheath edge; a subscript $\mathrm{w}$
denotes the wall. The symbol $f_{\mathrm{i}}(\bm{v})$ denotes the ion velocity
distribution function; angle brackets $\langle\cdot\rangle$ denote a
velocity-space average. Reconstructed (inverted) quantities carry a hat, e.g.
$\hat{\Gamma}_{E}$, to distinguish them from the underlying true value
$\Gamma_{E}$.
